% !TEX program = lualatex
% 日本語モード（修士論文）のサンプル
% ビルド: latexmk sample-ja.tex   （latex-devkit なら make build-local MAIN=sample-ja.tex）
%
% オプション:
%   ja / en            本文の言語（既定 ja）
%   master / doctor    修士論文 / 博士論文（既定 master）
%   print              両面印刷用（綴じ代つき・章は奇数ページから・リンクは黒）
%   draft              下書き（行番号・各ページにバージョンとビルド日時）
%   showyear           表紙に年度を表示（年度はビルド日から自動計算）
%   officialcolor      表紙の帯を技大 VI の青 #00A4E5 そのままにする
%   10pt / 11pt / 12pt 本文の文字サイズ（既定 11pt）
\documentclass[ja,master]{tmlthesis}

% 表紙の色を変えるとき（16進のRGB）
% \tmlthesissetup{maincolor=2F3E46, accentcolor=84A98C}

\jtitle{誤り訂正出力符号に基づく多値分類法の\\性能評価に関する研究}
\etitle{Performance Evaluation of Multiclass Classifiers\\Based on Error-Correcting Output Codes}
\jauthor{長岡 太郎}
\eauthor{Taro NAGAOKA}
\studentid{24000000}
\jsupervisor{雲居 玄道 准教授}
\jfield{情報・経営システム工学分野}
\efield{Information and Management Systems Engineering}

% \draftversion{v0.3}   % draft オプション時に表示するバージョン

\begin{document}

\frontmatter
\maketitle

\begin{abstract}
This thesis studies multiclass classification based on error-correcting
output codes (ECOC). We analyze the classification error of ECOC
classifiers and show conditions under which the code-word table attains
the optimal error rate. Experiments on text classification support the
analysis.
\end{abstract}

% 和文の概要を本文に入れる場合
% \begin{jabstract}
% 本論文では……
% \end{jabstract}

\tableofcontents
\listoffigures
\listoftables

\mainmatter
\chapter{序論}

\section{研究の背景}

特徴ベクトルからデータの所属カテゴリを推定する問題を分類問題という。
カテゴリ数 $M \geq 3$ の多値分類問題に対しては、複数の2値分類器を組み合わせる
誤り訂正出力符号（Error-Correcting Output Codes; ECOC）\cite{dietterich1994}が知られている。
2値分類器としては Support Vector Machine \cite{vapnik1998} などが用いられる。
ECOC の最適性については文献\cite{kumoi2021}で議論されている。

和文と欧文が混在する箇所の強調も確認しておく。
\textbf{太字の和文と Bold Latin}、\textit{斜体の欧文 italic text と和文}、
\textsf{ゴシック体と sans-serif} はいずれも組める\footnote{%
  ltjs 系クラスでは luatexja-preset の deluxe 指定がないと、
  和欧混在の斜体でフォントエラーになることがある。}。

\section{本論文の構成}

\begin{itemize}
  \item 第\ref{chap:method}章では、提案手法と理論的な性質を述べる。
  \item 第\ref{chap:conclusion}章では、結論と今後の課題を述べる。
\end{itemize}

\chapter{提案手法}\label{chap:method}

\section{定式化}

カテゴリ集合を $\mathcal{Y} = \{1, \dots, M\}$、符号語表を
$\bm{W} \in \{0,1\}^{M \times L}$ とする。入力 $\bm{x}$ に対する推定カテゴリは
\begin{equation}
  \hat{y} = \operatorname*{argmin}_{y \in \mathcal{Y}} \sum_{l=1}^{L} \bigl| w_{yl} - f_l(\bm{x}) \bigr|
  \label{eq:decode}
\end{equation}
で与えられる。

\begin{definition}
符号語表の任意の2行の間のハミング距離の最小値を、最小距離 $d_{\min}$ という。
\end{definition}

\begin{theorem}\label{thm:bound}
各2値分類器の誤り率が $p < 1/2$ で独立ならば、式\eqref{eq:decode}の誤り率は
$d_{\min}$ について単調に減少する。
\end{theorem}

\begin{proof}
2値分類器の誤り数が $\lfloor (d_{\min}-1)/2 \rfloor$ 以下なら正しく復号される。
この事象の確率は $d_{\min}$ について単調に増加する。
\end{proof}

\section{実験}

図\ref{fig:pipeline}に処理の流れを、表\ref{tab:result}に結果を示す。

\begin{figure}[tb]
  \centering
  \begin{tikzpicture}[box/.style={draw,rounded corners,minimum height=8mm,minimum width=22mm}]
    \node[box] (a) {入力 $\bm{x}$};
    \node[box,right=10mm of a] (b) {2値分類器 $f_l$};
    \node[box,right=10mm of b] (c) {復号};
    \draw[->] (a) -- (b);
    \draw[->] (b) -- (c);
  \end{tikzpicture}
  \caption{ECOC による分類の流れ}\label{fig:pipeline}
\end{figure}

\begin{table}[tb]
  \centering
  \caption{分類精度の比較}\label{tab:result}
  \begin{tabular}{lcc}
    \toprule
    手法 & 正解率 & 符号長 $L$ \\
    \midrule
    1-vs-Rest & 0.812 & 10 \\
    提案法    & \textbf{0.846} & 15 \\
    \bottomrule
  \end{tabular}
\end{table}

\chapter{結論}\label{chap:conclusion}

本論文では、ECOC に基づく多値分類法の誤り率を解析し、
定理\ref{thm:bound}により最小距離と誤り率の関係を示した。
今後の課題は、2値分類器の誤りが独立でない場合への拡張である。


\backmatter
\begin{acknowledgements}
本研究を進めるにあたり、熱心にご指導いただいた雲居玄道准教授に深く感謝いたします。
\end{acknowledgements}

\bibliographystyle{jIEEEtran}
\bibliography{references}

\end{document}
